\par\medskip
\noindent\textbf{Corolario 6.3.}
Sea $X_0$ un $k$-esquema de tipo finito y sea $x\in X_0$ una
singularidad cuadrática no degenerada racional sobre $k$.

\smallskip
\noindent\textbf{(1)} Sean $(R,\mathfrak X)$ y
$(R,\mathfrak Y)$ dos deformaciones de $X_0$, con
$R\in\operatorname{ob}(\widehat{\underline C}_\Lambda)$. Supongamos
que $k=k^2$. Entonces
\[
\mathfrak X_{(\widehat x)}\simeq
\mathfrak Y_{(\widehat x)}
\]
como esquemas formales $R$-ádicos si y solo si
\[
I_R(\mathfrak X_{(x)}|R)=I_R(\mathfrak Y_{(x)}|R).
\]

\smallskip
\noindent\textbf{(2)} Sea
$R\in\operatorname{ob}(\widehat{\underline C}_\Lambda)$ un anillo
local regular y sea $(R,\mathfrak X)$ una deformación de $X_0$.
Entonces $\mathcal O_{\mathfrak X,x}$ es regular si y solo si
\[
I_R(\mathfrak X_{(x)}|R)\not\subset\mathfrak m_R^2;
\]
equivalentemente, $I_R(\mathfrak X_{(x)}|R)$ es un divisor regular en
$\operatorname{Spf}(R)$.

\smallskip
\noindent\textit{Demostración.}
Sea $(R_0,Z)$ una deformación versal formal de
$\operatorname{Spec}(\widehat{\mathcal O}_{X_0,x})$. Como $x$ es una
singularidad cuadrática no degenerada racional sobre $k$, el Lema 6.2
nos permite tomar
\[
R_0\simeq\Lambda[[t]],
\qquad
Z\simeq\operatorname{Spf}
 \bigl(R_0[[X_1,\ldots,X_n]]/(F-t)\bigr).
\]

Para (1), las deformaciones
$(R,\mathfrak X_{(\widehat x)})$ y
$(R,\mathfrak Y_{(\widehat x)})$ son inducidas por $(R_0,Z)$ mediante
homomorfismos locales continuos de $\Lambda$-álgebras
\[
\varphi,\psi:R_0\longrightarrow R.
\]
Pongamos $a=\varphi(t)$ y $b=\psi(t)$. Entonces
\[
\begin{aligned}
\mathfrak X_{(\widehat x)}
&\simeq\operatorname{Spf}
 \bigl(R[[X_1,\ldots,X_n]]/(F-a)\bigr),\\
\mathfrak Y_{(\widehat x)}
&\simeq\operatorname{Spf}
 \bigl(R[[X_1,\ldots,X_n]]/(F-b)\bigr).
\end{aligned}
\]

%%%%%%%%%%  scan idx 138  |  folio 127  %%%%%%%%%%

Supongamos que
$I_R(\mathfrak X_{(x)}|R)=I_R(\mathfrak Y_{(x)}|R)$;
equivalentemente, $(a)=(b)$ en $R$. Entonces $a=ub$ para alguna
unidad $u\in R$. Como $k=k^2$ y $R$ es completo, $u=v^2$ para alguna
unidad $v\in R$, y por tanto $a=v^2b$. Así,
\[
(F-a)\longmapsto v^2(F-b)
\]
bajo $X_i\mapsto vX_i$, y esta sustitución induce un isomorfismo de
$R$-álgebras
\[
R[[X_1,\ldots,X_n]]/(F-a)
 \xrightarrow{\ \sim\ }
R[[X_1,\ldots,X_n]]/(F-b),
\]
porque $F$ es cuadrática. El recíproco es inmediato.

Para (2), existe un homomorfismo local continuo de $\Lambda$-álgebras
\[
\varphi:R_0=\Lambda[[t]]\longrightarrow R
\]
tal que
\[
\mathfrak X_{(\widehat x)}
\simeq\operatorname{Spf}
 \bigl(R[[X_1,\ldots,X_n]]/(F-\varphi(t))\bigr).
\]
Además,
\[
I_R(\mathfrak X_{(x)}|R)
=I_R(\mathfrak X_{(\widehat x)}|R)
=(\varphi(t)),
\]
y
\[
R[[X_1,\ldots,X_n]]/(F-\varphi(t))
=\widehat{\mathcal O}_{\mathfrak X,x}
\]
es regular si y solo si $\mathcal O_{\mathfrak X,x}$ es regular. Como
$R$ es regular, $R[[X_1,\ldots,X_n]]$ es un anillo local regular. Si
$M$ denota su ideal maximal, entonces su cociente por
$F-\varphi(t)$ es regular si y solo si
\[
F-\varphi(t)\notin M^2,
\]
lo cual equivale a $\varphi(t)\notin\mathfrak m_R^2$, y por tanto a
\[
I_R(\mathfrak X_{(x)}|R)\not\subset\mathfrak m_R^2.
\]
\hfill$\square$
